\chapter{Bulimia Nervosa}

\section*{Definition and Overview}
Bulimia nervosa is a serious eating disorder characterised by repeated episodes of binge eating, in which unusually large amounts of food are consumed in a short time with a sense of loss of control, followed by compensatory behaviours designed to undo the effects of the binge. These behaviours include self-induced vomiting, misuse of laxatives or diuretics, fasting, or excessive exercise. Because people with bulimia are often of normal or slightly above-normal weight, the disorder can remain hidden for years. It is driven by an intense fear of gaining weight and an overvaluation of body shape and weight, and it carries significant medical and psychological risks that require coordinated treatment.

\section*{Epidemiology}
Bulimia nervosa affects roughly 1--2\% of women at some point in their lives and is approximately ten times more common in women than in men. It usually begins in late adolescence or early adulthood, most often between the ages of 16 and 25, and it occurs across all cultural and socioeconomic backgrounds, though it is more common in societies that emphasise thinness. Only a minority of people with bulimia are ever diagnosed, because the bingeing and purging are usually kept secret and the individual often appears healthy. The disorder carries a substantial burden of co-existing depression, anxiety, and substance use.

\section*{Aetiology and Pathophysiology}
Bulimia develops from an interaction of genetic, psychological, sociocultural, and behavioural factors.

\begin{itemize}
    \item \textbf{Genetic factors:} a family history of eating disorders increases risk, and twin studies show a substantial heritable component
    \item \textbf{Psychological factors:} low self-esteem, perfectionism, and impulsivity are common features
    \item \textbf{Dieting:} most cases begin after a period of strict dietary restriction, which then gives way to disinhibited binge eating
    \item \textbf{Cultural pressure:} social and media emphasis on thinness and body image contributes to dissatisfaction with weight and shape
    \item \textbf{Life stress:} difficult events, conflict, or emotional distress commonly trigger bingeing as a coping response
    \item \textbf{Physiological disruption:} repeated bingeing and purging disturb hunger signals, satiety, and electrolyte balance, reinforcing the cycle
\end{itemize}

\section*{Clinical Features}
The disorder is marked by a recurring cycle of bingeing and purging, which may be accompanied by the following features.

\begin{itemize}
    \item Recurrent episodes of eating unusually large amounts of food with a subjective sense of loss of control
    \item Compensatory purging by self-induced vomiting, laxatives, diuretics, fasting, or excessive exercise
    \item Intense fear of weight gain and harsh self-criticism about body shape and size
    \item Eating in secret, with feelings of shame and guilt afterwards
    \item Swelling of the parotid glands (cheeks and jaw) and erosion of the back surfaces of the teeth from repeated exposure to stomach acid
    \item Calluses or scarring on the knuckles (Russell's sign) from inducing vomiting
    \item Sore throat, heartburn, and irregular or absent menstrual periods
    \item Mood swings, anxiety, and avoidance of social eating
\end{itemize}

\section*{Differential Diagnosis}
Similar eating problems and other conditions should be considered.

\begin{itemize}
    \item \textbf{Anorexia nervosa, binge-purge type:} bingeing and purging occurring in the presence of significantly low body weight
    \item \textbf{Binge eating disorder:} binge eating without compensatory behaviours
    \item \textbf{Other specified feeding or eating disorders:} atypical patterns that do not meet full criteria
    \item \textbf{Depression and anxiety disorders:} commonly co-exist and may drive or be driven by the eating disorder
    \item \textbf{Body dysmorphic disorder:} preoccupation with a perceived flaw in appearance
    \item \textbf{Gastrointestinal and endocrine disorders:} conditions causing vomiting or metabolic disturbance that mimic purging effects
\end{itemize}

\section*{When to Seek Medical Care}
See a doctor if you recognise bingeing and purging in yourself or in someone else. Bulimia responds best to early treatment, and general practitioners, psychologists, and dietitians are experienced in helping people recover.

\textbf{Call 000 (or local emergency services) or seek urgent help for:}
\begin{itemize}
    \item Thoughts of suicide or self-harm -- call 000 or Lifeline on 13 11 14
    \item Vomiting blood or severe abdominal pain
    \item Fainting, dizziness, or a very rapid or irregular heartbeat
    \item Swelling of the hands and feet, or passing very little urine
    \item Severe dehydration from repeated vomiting or laxative misuse
\end{itemize}

\section*{Diagnosis and Investigations}
The diagnosis is made through a careful clinical interview, supported by physical assessment for the effects of the disorder.

\begin{itemize}
    \item \textbf{Clinical interview:} a sensitive discussion of eating habits, bingeing, purging, weight, and body image, using screening tools such as the SCOFF questionnaire where appropriate
    \item \textbf{Physical examination:} assessment of weight, dental enamel, parotid glands, hands, throat, and signs of dehydration
    \item \textbf{Blood tests:} electrolytes, particularly potassium, plus renal function and other markers of purging
    \item \textbf{Electrocardiogram:} to detect cardiac rhythm disturbance from hypokalaemia
    \item \textbf{Bone density scanning:} in some cases, because nutritional deficiency can weaken bones
    \item \textbf{Mental health assessment:} screening for depression, anxiety, and risk of self-harm
\end{itemize}

\section*{Management}
Treatment is usually delivered on an outpatient basis and addresses both the eating behaviour and its psychological drivers.

\begin{itemize}
    \item \textbf{Cognitive behaviour therapy (CBT):} the most effective approach, focused on breaking the binge--purge cycle and changing unhelpful beliefs about food, weight, and shape
    \item \textbf{Family-based treatment:} involving parents and family members, particularly for adolescents and younger people
    \item \textbf{Dietary counselling:} structured, regular meals planned with a dietitian to restore normal eating patterns
    \item \textbf{Psychological support:} treatment of co-existing anxiety, depression, or low mood
    \item \textbf{Medication:} the antidepressant fluoxetine, at a higher dose than used for depression, reduces bingeing in some people
    \item \textbf{Dental care:} a dentist to manage the damage caused by gastric acid on the teeth
    \item \textbf{Hospital care:} required for severe complications such as marked electrolyte disturbance, dehydration, or suicidal intent
\end{itemize}

\section*{Complications}
Repeated bingeing and purging can damage many body systems.

\begin{itemize}
    \item Dental erosion, decay, and throat irritation from repeated exposure to stomach acid
    \item Low potassium and other electrolyte imbalances, which can cause dangerous cardiac arrhythmias
    \item Dehydration and kidney damage
    \item Enlarged salivary glands and a chronically sore throat
    \item Inflammation, tearing, or rare rupture of the stomach or oesophagus
    \item Irregular periods, fertility problems, and osteoporosis
    \item Laxative dependence and chronic bowel dysfunction
    \item Co-existing depression and anxiety, with an elevated risk of suicide
\end{itemize}

\section*{Prognosis and Outcomes}
With treatment, roughly half of people with bulimia nervosa recover fully, and many others improve substantially. Recovery often takes months to years, and relapses are common during periods of stress, but the long-term outlook is more favourable than for anorexia nervosa. Early treatment and strong family support improve the chances of recovery, and even people who do not fully recover usually learn to manage their eating and live fulfilling lives.

\section*{Prevention and Self-Care}
Prevention at a community level involves promoting healthy body image, but individuals and families can also take practical steps.

\begin{itemize}
    \item Challenge unrealistic ideas about weight and shape, and resist diet-culture pressures
    \item Eat regular, balanced meals rather than skipping food, which can trigger bingeing
    \item Avoid strict dieting, and use professional support if weight change is genuinely needed
    \item Seek help early, since a shorter duration of illness is linked to better recovery
    \item Develop non-food ways of coping with stress and difficult emotions
    \item Encourage open conversation about body image with children and teenagers
    \item Stay engaged with treatment and support services, and attend follow-up appointments
\end{itemize}

\section*{Sources and Further Reading}
The following organisations provide reliable information on bulimia nervosa and eating disorders for patients, families, and clinicians.

\begin{itemize}
    \item NHS. \textit{Bulimia nervosa: symptoms, treatment, and getting help.} https://www.nhs.uk/mental-health/conditions/bulimia/
    \item Mayo Clinic. \textit{Bulimia nervosa: symptoms, causes, and treatment.} https://www.mayoclinic.org/
    \item National Eating Disorders Collaboration. \textit{Eating disorders information and resources for families and professionals.} https://www.nedc.com.au/
    \item MedlinePlus. \textit{Bulimia nervosa health information.} https://medlineplus.gov/bulimia.html
    \item National Institute for Health and Care Excellence. \textit{Eating disorders: recognition and treatment guidance.} https://www.nice.org.uk/
    \item Butterfly Foundation. \textit{Eating disorder support, counselling, and helpline services.} https://www.butterfly.org.au/
\end{itemize}
